% Abstract and index terms (issue #7).

\begin{abstract}
Remaining useful life (RUL) prediction estimates how many operating cycles an
engine has left before it fails. This report presents Phase~1 of a project that
compares engineered degradation features with sequence models on the NASA C-MAPSS
turbofan data: feature-based machine-learning baselines on subset FD001, under a
protocol that Phase~2 will reuse. Training engines are split by engine, every
preprocessing step is refitted inside each training fold, the RUL target is
uncapped, and each of the 100 official test engines is scored once at its last
cycle. From regime-normalized sensor readings we build 155 causal predictors
(rolling means, slopes and changes since each engine's first cycle) and tune
Random Forest, XGBoost and Elastic Net regressors on them and on 25 raw inputs.
The engineered Random Forest reaches a test RMSE of 27.55 cycles and an MAE of
17.67 cycles, and all six models lie within five RMSE cycles of each other. The
engineered features roughly halve the tree models' error near failure (RMSE 7.52
and 5.62 against 13.50 and 14.32 with raw inputs) but overestimate young engines,
and a single engine contributes 92\% of the Random Forest's NASA score. Capping
the training target at 125 cycles lowers the Random Forest's RMSE to 16.21 with
engineered and 18.66 with raw inputs, reversing their order.
\end{abstract}

\begin{IEEEkeywords}
remaining useful life, predictive maintenance, C-MAPSS, feature engineering,
random forest, XGBoost, elastic net
\end{IEEEkeywords}
